\chapter{}
\label{app:spectrometric_system_code}

В данном приложении приведены структурная схема и листинг кода программы, использованной для обработки спектральных данных, формирования признакового описания и последующего анализа качества сварного соединения.

На рисунке~\ref{fig:spectra_algorithm_scheme} представлена общая схема программы. В листингах~\ref{lst:spectra_parser}--\ref{lst:spectra_analyze_results} приведены исходные коды основных модулей программы: парсера исходных файлов, класса спектра, модуля вычисления признаков и модуля анализа результатов.

\begin{figure}[H]
\centering
\begin{adjustbox}{max width=\linewidth}
\begin{tikzpicture}[
    node distance=0.4cm,
    every node/.style={font=\footnotesize},
    block/.style={
        rectangle,
        draw,
        rounded corners,
        align=center,
        minimum height=0.60cm,
        inner sep=2pt
    },
    arrow/.style={->, thick}
]

\node[block] (input) {Загрузка Raw8-файлов спектров и Excel-файла с дефектами};
\node[block, below=of input] (defects) {Извлечение признака $h_i$ и формирование таблицы дефектов};
\node[block, below=of defects] (raw) {Чтение Raw8-файлов канала 1712307U3: длины волн и интенсивности};
\node[block, below=of raw] (match) {Сопоставление спектров и дефектов по \texttt{sample\_id}};
\node[block, below=of match] (filterq) {Фильтрация шумовых спектров};
\node[block, below=of filterq] (features) {Расчёт признаков спектра: статистики, диапазоны длин волн, FFT};
\node[block, below=of features] (processing) {Дополнительная обработка: сглаживание, Savitzky--Golay, baseline};
\node[block, below=of processing] (table) {Формирование таблицы признаков \texttt{features.xlsx}};
\node[block, below=of table] (aggregate) {Агрегация признаков по образцу: среднее, СКО, минимум, максимум};
\node[block, below=of aggregate] (sampletable) {Формирование таблицы \texttt{sample\_features.xlsx}};
\node[block, below=of sampletable] (models) {Обучение и сравнение моделей: LogisticRegression, RandomForest, GradientBoosting, SVM, kNN};
\node[block, below=of models] (cv) {Оценка качества через StratifiedGroupKFold с группировкой по \texttt{sample\_id}};
\node[block, below=of cv] (sensitive) {Подбор чувствительного режима со сдвигом порога};
\node[block, below=of sensitive] (metrics) {Расчёт метрик: accuracy, balanced accuracy, precision, recall, F1, confusion matrix};
\node[block, below=of metrics] (save) {Сохранение таблиц, графиков, матрицы ошибок и важности признаков};
\node[block, below=of save] (end) {Интерпретация результатов и вывод о применимости спектральных признаков};

\draw[arrow] (input) -- (defects);
\draw[arrow] (defects) -- (raw);
\draw[arrow] (raw) -- (match);
\draw[arrow] (match) -- (filterq);
\draw[arrow] (filterq) -- (features);
\draw[arrow] (features) -- (processing);
\draw[arrow] (processing) -- (table);
\draw[arrow] (table) -- (aggregate);
\draw[arrow] (aggregate) -- (sampletable);
\draw[arrow] (sampletable) -- (models);
\draw[arrow] (models) -- (cv);
\draw[arrow] (cv) -- (sensitive);
\draw[arrow] (sensitive) -- (metrics);
\draw[arrow] (metrics) -- (save);
\draw[arrow] (save) -- (end);

\end{tikzpicture}
\end{adjustbox}
\caption{Общая схема алгоритма обработки спектральных данных и классификации качества сварки}
\label{fig:spectra_algorithm_scheme}
\end{figure}

Листинг~\ref{lst:spectra_parser} содержит модуль \texttt{Parser.py}, реализующий чтение исходных файлов спектральных измерений и извлечение данных, необходимых для последующей обработки.

\lstinputlisting[
    language=Python,
    caption={Модуль \texttt{Parser.py} для чтения и первичной обработки спектральных данных},
    label={lst:spectra_parser}
]{code/spectra/Parser.py}

Листинг~\ref{lst:spectra_spectrum} содержит модуль \texttt{Spectrum.py}, в котором реализовано представление спектра и связанных с ним параметров образца.

\lstinputlisting[
    language=Python,
    caption={Модуль \texttt{Spectrum.py} для представления спектральных данных},
    label={lst:spectra_spectrum}
]{code/spectra/Spectrum.py}

Листинг~\ref{lst:spectra_make_features} содержит модуль \texttt{make\_features.py}, предназначенный для вычисления спектральных признаков и формирования таблиц признакового описания.

\lstinputlisting[
    language=Python,
    caption={Модуль \texttt{make\_features.py} для формирования признаков спектральных сигналов},
    label={lst:spectra_make_features}
]{code/spectra/make_features.py}

Листинг~\ref{lst:spectra_analyze_results} содержит модуль \texttt{analyze\_results.py}, в котором реализованы загрузка таблиц признаков, построение моделей классификации, оценка их качества и сохранение результатов анализа.

\lstinputlisting[
    language=Python,
    caption={Модуль \texttt{analyze\_results.py} для анализа результатов и сравнения классификаторов},
    label={lst:spectra_analyze_results}
]{code/spectra/analyze_results.py}

\endinput